\honest{The Meditation on Curiosity}{Eliezer Yudkowsky}{2007}

I ask whether rationalists are obliged to criticize their own beliefs. My answer is that if you do it from a sense of duty, you want to have investigated, and that is different from wanting to investigate. Roger Zelazny distinguished wanting to be an author from wanting to write, and Mark Twain said a classic is something everyone wants to have read and no one wants to read.

Inquiry done from duty stops early. You consider an objection, find a counterargument, and stop, once you feel you have done enough. You end up confident enough to keep your old plans, and not so confident that you, or other rationalists, would find it hard to believe. I give this failure two names, ``motivated stopping'' and ``purchase of rationalist satisfaction.'' I give it no example.

A curious person looks where belief is most likely to change, or at inquiries least like those already tried, and does not mind which way it changes; afterwards the beliefs should not look as they did at the start. The dutiful person stays on familiar ground, to get the inquiry over with and restore the old balance. I give no example of anyone being curious. For that I refer you to ``A Fable of Science and Politics,'' in another book, whose last character, Ferris, shows innocent curiosity: lightness and an eager reaching for evidence. Ursula Le Guin wrote that innocence has no strength against evil, but has strength for good; a rationalist's training is a risk worth taking, but we can try to keep that lightness. I quote myself, from ``The Twelve Virtues of Rationality'': if you believe you already know, or do not wish to know, your questioning will be purposeless.

Then I admit the problem. You cannot make yourself curious by willing it, any more than you can will a cold foot to feel warm. I concede: ``Sometimes, all we have is our mere solemn vows.'' So what can you do with duty?

My remedies are these. Watch for sparks of real interest in a dutiful investigation, and for painful possibilities you flinch from. Meditate on conservation of expected evidence. Recite a litany. Imagine the advantages of believing the answer you fear.

For the second remedy, I state the law correctly: your expected belief after seeing new evidence should equal your belief now. Then I go further: your belief should always be ``evenly poised to shift in either direction,'' and each new point should have ``equal potential to shift belief upward or downward'': at 70\%, you should be as ready to drop to 69\% as to rise to 71\%. \nb{These two phrasings get the mathematics wrong. The law says only that the average of your possible future beliefs, weighted by how likely each is, equals your present belief. A likely small rise can balance an unlikely large fall. Yudkowsky had said exactly this in August, in ``Conservation of Expected Evidence.''} If the argument you are looking at is not new, why look there? Is it where a curious person would look? Are you ``criticizing your belief at its strong points, rather than its weak points,'' or rehearsing evidence?

If you avoid rehearsing known support and lower your belief ``one tiny bite at a time'' from new evidence, you may manage to give it up entirely. \nb{Also in August, in ``The Importance of Saying Oops,'' Yudkowsky had called ``grudgingly conceding each millimeter of ground'' the mistake, and said it was ``better by far to see it all at once.''}

The Litany of Tarski replaces thoughts about the wonderful consequences of a diamond in a locked box: if the box contains a diamond, I want to believe that it does; if it does not, I want to believe that it does not; let me not become attached to beliefs I may not want. Then think about the possibility that there is no diamond, and the advantage of believing so if it is true. See also Gendlin's litany.

Last: if you find the slightest shred of real uncertainty, you should guard it ``like a forester nursing a campfire,'' since curiosity makes you light and eager and gives your questioning purpose.
